% Generated draft: requires mathematical and source review.
\documentclass[11pt]{article}
\usepackage[margin=25mm]{geometry}
\usepackage{fontspec}
\setmainfont{Latin Modern Roman}
\usepackage{amsmath,amssymb,mathtools,unicode-math}
\setmathfont{Latin Modern Math}
\usepackage{enumitem,needspace,xurl,hyperref}
\setlength{\parindent}{0pt}
\setlength{\parskip}{5pt}
\setlength{\emergencystretch}{3em}
\newcommand{\N}{\mathbb{N}}
\newcommand{\Z}{\mathbb{Z}}
\newcommand{\Q}{\mathbb{Q}}
\newcommand{\R}{\mathbb{R}}
\newcommand{\C}{\mathbb{C}}
\newcommand{\F}{\mathbb{F}}
\newcommand{\A}{\mathbb{A}}
\newcommand{\PP}{\mathbb{P}}
\newcommand{\E}{\mathbb{E}}
\newcommand{\eps}{\varepsilon}
\newcommand{\dd}{\,\mathrm{d}}
\begin{document}
\section*{UnsolvedMath ID: 11000089}
\textbf{Source problem code:} catalogue entry
\subsection{Definitions and mathematical statement}
Retain the space written \(T_g^c\) in the source's compressed notation ``T c g.'' Its definition and the meaning of the superscript \(c\) are absent, so the draft does not identify it with a compact core, compact-type Torelli space, or another completion without further evidence.

\emph{Problem 4.3.} Determine good upper and lower bounds for the homological dimension, or the CW dimension, of this space. In a fixed ordinary-homology convention with coefficients in \(R\), one possible invariant is
\[
 \operatorname{hd}_R(T_g^c)=\sup\{n:H_n(T_g^c;R)\neq0\}.
\]
A common homotopy-CW convention would instead use
\[
 \inf\{\dim X:X\text{ is a CW complex homotopy equivalent to }T_g^c\}.
\]
The excerpt does not fix coefficients, local-coefficient variants, actual versus homotopy CW dimension, or whether these are the intended definitions. It requests useful bounds rather than providing a proposed formula or a required sharp asymptotic order. The genus range and allowed dependence of bounds must also be specified.

The supplied research summary mentions \(\operatorname{cd}(\mathcal I_g)=3g-5\), attributing it to Bestvina--Bux--Margalit, and calls the target a compact core. The space and a comparison with a classifying space of \(\mathcal I_g\) are not supplied. A group cohomological-dimension result does not automatically determine ordinary homological dimension or CW dimension of this undefined \(T_g^c\). The quoted formula and compact-core interpretation remain source-summary claims requiring review, not assumptions used to derive bounds.
\subsection{Short English statement}
Bound the homological or CW dimension of the source's space \(T_g^c\). Its definition and dimension convention are missing, and the cited Torelli-group cohomological-dimension result cannot be transferred without a justified comparison.
\subsection{Sources}
\begin{itemize}
\item Local UnsolvedMath registry, permanent ID 11000089, source code catalogue entry, Problem 4.3, PDF page 76. Repeated problem numbering is distinguished by this page. No individual proposer is supplied.
\item Farb, \emph{Problems on Mapping Class Groups and Related Topics}. Supplied URL: \url{https://www.math.uchicago.edu/~farb/papers/mcgbook.pdf}; recorded access 2026-07-29; extraction pdf-font-structure. The edited volume contains 329 explicitly typeset headings across contributed chapters.
\item The supplied summary cites Bestvina--Bux--Margalit for \(\operatorname{cd}(\mathcal I_g)=3g-5\) and mentions partial bounds on a compact core, without full bibliography, genus hypotheses, or a comparison theorem identifying the target space.
\item Supplied status open and OPEN-TRIAGE, with editorial partial-progress language. Current primary-literature status and redistribution terms for the publicly accessible author/publisher copy remain NEEDS\_REVIEW.
\end{itemize}
\end{document}
